\externaldocument{./parts/project/chapters/TransportSystemApi.tex}
\externaldocument{./parts/project/chapters/TransportSystemImpl.tex}

\chapter{General Development Details} \label{planningTheProject}
\section{Development in C}
Since a key desire when designing the initial task for
the master thesis was to reach as many developers as possible,
especially lower-level developers who normally work with sockets,
it was decided that the inital implementation should be done in C.
C forms the backbone of many modern programming libraries, and lack
many of the quality of life networking features of higher level languages.
By implementing TAPS in C, it would also be possible to try to integrate it
into a software library such as glib. \todo{What library did michael actually talk about here?}

Integrating a TAPS implementation into a larger library would be the easiest
way of making it available to more developers, as they would not have to
rely on a relatively unknown implementation, in order to gain the benefits
of transport services.

Apart from the ubiquity, C offers several advantages. First and foremost
is the raw performance of C and the maturity of software libraries written
for it. If you are for example writing a high performance networking server,
then C or C++ is a natural choice. The maturity of many available software libraries is
another big advantage with C. For example, since TAPS is designed to be asynchronous,
it needed to run an event loop as a core part of the application. For this,
several mature libraries, which are already used in large applications are
available, such as libuv, which is used to form the core event loop of NodeJS. \cite{Libuv2025}

\subsection{Downsides of C}
The abstract interface described by the TAPS WG is very clearly object
oriented, C is a procedural programming language, and does not feature
the concept of classes or objects, which would most neatly map to the
TAPS interface. Instead, the object orientation and division of responsibility
featured in an object-oriented language has to be emulated through the use of
structs, and dedicated functions which operate on them. An example of this
would be the struct representing remote endpoints, which contains all the
context related to remote endpoints, which is for example used
in Connections and Listeners:

\begin{figure}[H]
\begin{minted}{C}
typedef struct RemoteEndpoint{
  RemoteEndpointType type;
  uint16_t port;
  union {
    struct sockaddr_storage address;
    char* hostname;
  } data;
} RemoteEndpoint;
\end{minted}
\caption{Remote endpoint type in C}
\end{figure}

Along with the definition of this struct, the files containing the
source code for remote endpoints also contain functions to create 
remote endpoints with different types of data, in ways which might be useful
to an application programmer, for example for initially clearing the struct
and building it with an IPv4 address:

\begin{figure}[H]
\begin{minted}{C}
void remote_endpoint_build(RemoteEndpoint* remote_endpoint) {
  remote_endpoint->data.address.ss_family = AF_UNSPEC;
  memset(&remote_endpoint->data, 0, sizeof(remote_endpoint->data));
}

void remote_endpoint_with_ipv4(
    RemoteEndpoint* remote_endpoint, 
    in_addr_t ipv4_addr) {
  remote_endpoint->type = REMOTE_ENDPOINT_TYPE_ADDRESS;

  struct sockaddr_in* addr = (struct sockaddr_in*)&remote_endpoint->data.address;
  addr->sin_family = AF_INET;
  addr->sin_addr.s_addr = ipv4_addr;
  addr->sin_port = htons(remote_endpoint->port);
}
\end{minted} 
\caption{Remote endpoint utility functions}
\end{figure}
\todo{Update this when the final code is done}

By defining structs and dedicated functions to work on the structs like
above, we can to a very large degree emulate the encapsulation of responsibilites
as seen in a object oriented language. By encapsulating the data structures in the
above way, we can also ensure that the external API which is exposed to the end user,
is as similar as possible to the API described in RFC9622 \cite{trammellAbstractApplicationProgramming2025}

The above code snippets also show another common design pattern used in C, which
is the use of output parameters instead of return values. In C there is no garbage
collector, and no built in support for reference counting, such as with \texttt{std::shared\_ptr}
in C++ for example, the programmer has to manage allocation of memory themselves.
If we could not use output parameters, the function would have to allocate memory on
the heap itself, before returning a pointer to this region of memory, expecting the
user of the library to later free this region of memory. 
This will lead to memory leaks
if the user of the library is not aware that they have to do this. 
The function cannot
return a pointer to a variable it has defined on the stack, as this stack
region can be overwritten after the function has returned, and will lead to a dangling
pointer.
By using output parameters as often as possible, the library avoids having to
allocate memory other than when strictly necessary. If the user of a library
wants a long-lived struct, they can allocate memory on the heap themselves, before
passing a pointer to this region of memory to a library function. This has
the added benefit of making the library user more aware of when memory is allocated,
minimizing the risk of them forgetting to free memory.

\section{Testing the project} \label{testingTheProject}
A core tenet since the very start of development was that the library should be
as robust as possible. This would to a very large degree achieved primarily through
the use of automated testing, and continious integration. By writing automated tests
which run on every commit to the repository, it would be easy to see exactly which change
has lead to a bug, making it far easier to fix. A well-tested library is also an absolute
necessity when trying to attract external dependents, or to integrate it into a larger software
library. This is especially important for large libraries such as TAPS, which abstract away
large amounts of logic away from the application programmer, making it harder to detect bugs
and know exactly where the problem in their application lies.

Writing tests for an asynchronous networking library is not necessarily straight-forward.
Both networking and asynchronous programming is not necessarily deterministic, which makes
it harder to write reliable tests which can verify the library's expected behaviour. Another
challenge is that there are not very many testing libraries available in pure C. Early on,
it was therefore decided to write the tests in C++, which allows for interoperability with the C
library, but would also allow for using a C++ testing library such as Google test. \cite{GoogleTest2025}
This allows for a richer feature set than a simpler library or manually written tests in C,
and allows for a richer set of assertions, verifying for example that the library fails in the expected
manner when given invalid data.

Tests are first separated into type and then into responsibility. When separating
by type, the main types are integration tests and unit tests. Integration tests perform a larger check on the
whole library and often actually interact with the network by sending packets or creating
connections to a local server. \todo{Talk about how this is done, maybe via docker?}
Unit tests perform a smaller check on individual functions, attempting to verify
all logic and that differing branches are hit when expected, based on the input data.
When dividing by responsibility, the testing stack largerly follows the same model
as the main library, with test files separated by responsibility of the original
file or struct that it is performing tests on.

A central concept when writing good tests is to isolate the component being tested from
external libraries. This is important for several reasons. First of all, we are not interested
in verifying the behavior of our external dependencies, as we have to assume that they
work well in order to limit the scope of our testing. Secondly, we need to know that if
an external library \textit{does} fail, we need to know that our library will handle it
in the correct manner. Therefore we have to simulate that an external library fails and
verify our applications behaviour in this case. All of these concerns are best met by
mocking. Mocking is the concept of defining a custom behaviour for an external library
or dependency when testing, instead of invoking the real library. By using this we can
simulate the full range of behaviours from the external library and verify that our application
reacts correctly, while also avoiding time-consuming API calls, making unit tests
significantly faster. For mocking, CTaps decided on FFF, which is a small open-source
framework for creating fake C functions when testing. \cite{longFakeFunctionFramework2025}
FFF can either mock the external library functions, such as \texttt{uv\_getaddrinfo},
which is the built in support for DNS lookup in libuv, or mock a call to a different
function, which either contains a call to an external library, or which we are simply
not interested in testing.

\begin{figure}
\begin{minted}{Cpp}
DEFINE_FFF_GLOBALS;
FAKE_VALUE_FUNC(int, uv_udp_bind, uv_udp_t*, const struct sockaddr*, unsigned int)
FAKE_VALUE_FUNC(int, uv_udp_recv_start, uv_udp_t*, uv_alloc_cb, uv_udp_recv_cb)

TEST(InitiationTests, respectsLocalEndpoint) {
    uv_udp_bind_fake.return_val = 0;
    uv_udp_recv_start_fake.return_val = 0;
    
    ... Rest of code which calls uv_udp_bind and uv_udp_recv_start ...
    
    // Assert that the mocked functions were invoked with the expected arguments    
    EXPECT_EQ(1, uv_udp_bind_fake.call_count);
    EXPECT_EQ(1234, htons(((struct sockaddr_in*)uv_udp_bind_fake.arg1_val)->sin_port));
    ... Rest of the assertions ...
}

\end{minted}
\caption{Example test mocking udp functions, using FFF}
\end{figure}


\section{Code Style}
Because CTaps was devised as an open source library, which
should either provide a usable solution for application
developers looking to make use of Transport Services, or as
a useful branching off point for a larger implementation, CTaps has aimed
to keep a consistent, opionated coding style, in order to keep a look and
feel which remains consistent, both with the TAPS API, and which remains
consistent within the project itself. This style was early on decided that
it should follow an established style guide, making as few changes to the configuration
as possible, in order to make the project seem as familiar as possible to external
users and developers. The chosen style guide was the Google C++ style guide, which
where the subset of rules applicable to C, would be applied to the project \cite{googleGoogleStyleGuide}

The formatting of CTaps is automated using clang-format, which can parse
and automatically fix simple format changes \todo{citation?}. This has built-in
support for the Google style guide. \todo{Write about verifying the formatting in github actions?}

\section{Development Methodology} \label{developmentMethodology}
Since this thesis was developed by one person, it has not
followed a strict development methodology such as Scrum, however,
it has tried to be as agile as possible, working on development of minimal features
at a time. A detailed overview of the development timeline and iterations can be
seen in chapter \ref{projectTimeline}. Instead of carefully planning how the TAPS
API concepts would be implemented, they were instead implemented in a a minimal state,
later revising as more features were needed. An example of this is Connections, which
first were implemented with only support for UDP, later updating them to be generic,
in order to support TCP and other protocols.

Developing in small iterations, ensuring correct functionality at each iteration, helps
limit the complexity, both when first implementing, but especially when trying to fix
issues with the implementation. By trying to ensure correctness at each step in the development,
one can to as a large degree as possible limit where bugs will arise. This ties
into another important factor in the development of this thesis, which is the use of
Test-Driven-Development (TDD). When developing a new feature, a test
verifying the functionality of this feature was developed alongside it.
This verifies that the new functionality is actually functional,
while ensuring that breaking changes are not introduced at a later stage.
A detailed overview of the test system can be seen in section \ref{testingTheProject}


\section{External dependencies}
\subsection{Dynamically Linked Dependencies?} \todo{Are these actually dynamically linked?}
\subsubsection{Libuv}
As mentioned in chapter \ref{planningTheProject}, CTaps is reliant on several
external libraries, both for testing and for core functionality. The largest
and most essential of these are libuv \cite{Libuv2025}. Libuv provides the central
event loop which allows CTaps to be asynchronous. It also provides several quality
of life features, such as built in support for asynchronous interactions with TCP
and UDP, through the \texttt{uv\_tcp\_t} and \texttt{uv\_udp\_t} handles respectively.
When adding protocols which libuv does not have native support for, such as QUIC,
it offers a custom work handle \texttt{uv\_work\_t}, which can be used to define
a custom piece of work. This can be used to define both a long-lived QUIC-based
server, or a short interaction with an existing server. Which underlying handle
is actually used in libuv is naturally abstracted away in CTaps, and for the
end-user of the library the interaction is the same no matter what. \todo {very informal lol}

\subsubsection{Glib}
In order to increase the reliability of internal data structures, as well as
minimizing the amount of time spent reinventing the wheel, CTaps relies on the
Glib library. Glib is a low level core library which is used at the core of
larger libraries such as GTK and GNOME \todo{citation}. Glib provides several
key data structures used internally in CTaps. Among others, it provides the
Hash Map used to keep track of received Connections in a Listener, the tree
used to gather candidates in the Candidate Gathering algorithm, among others.
While adding an external library, especially in core functionaly such as networking,
poses a risk, the benefits were deemed to be greater than the costs. Primarily,
this consideration was made on the basis of Glib being a relatively high-profile
library used in several large projects already, in addition to the benefits of
increased development speed, since CTaps would not have to implement its own 
basic data structures. Using a well-tested external library has the added benefit
of avoiding bugs, as if CTaps rolled its own data structures, it would have
to worry not only about bugs in the library logic, but also in the
data structures themselves.

In order to minimize how much a library consumers would have to
adhere to Glib, it was decided that the external API should be
kept free of Glib data structures, and that all external functions
should only take parameters defined in standard C or in CTaps itself.

\subsection{Statically Linked Dependencies}
Minor software which does not provide a core functionality to CTaps
is statically linked and is therefore distributed along with the
compiled executable. \todo{Hvordan blir biblioteket egentlig formatert}
This is done in order to make installation of
CTaps as easy as possible. Since the libraries are small, they
do not bloat the CTaps binary to the same degree as a larger library
such as libuv would, if it were to be statically linked.

\subsubsection{log.c}
Logging
https://github.com/rxi/log.c

\section{Build System}
Should we describe the build system? Automated testing which I will set up?
CMake? automated packaging etc. etc. 

\section{Architecture}
CTaps follows the overarching architecture as described in RFC 9621 \cite{paulyArchitectureRequirementsTransport2025}.
This splits a TAPS system into two two main parts, a Transport Services
API as described in \cref{transportSystemApi} and a Transport System Implementation
as described in \cref{transportSystemImpl}.

\begin{figure}[H]
    \centering
    \begin{lstlisting}
     +-----------------------------------------------------+
     |                    Application                      |
     +-+----------------+------^-------+--------^----------+
       |                |      |       |        |
     pre-               |     data     |      events
     establishment      |   transfer   |        |
       |        establishment  |   termination  |
       |                |      |       |        |
       |             +--v------v-------v+       |
     +-v-------------+   Connection(s)  +-------+----------+
     |  Transport    +--------+---------+                  |
     |  Services              |                            |
     |  API                   |  +-------------+           |
     +------------------------+--+  Framer(s)  |-----------+
                              |  +-------------+
     +------------------------|----------------------------+
     |  Transport             |                            |
     |  System                |        +-----------------+ |
     |  Implementation        |        |     Cached      | |
     |                        |        |      State      | |
     |  (Candidate Gathering) |        +-----------------+ |
     |                        |                            |
     |  (Candidate Racing)    |        +-----------------+ |
     |                        |        |     System      | |
     |                        |        |     Policy      | |
     |             +----------v-----+  +-----------------+ |
     |             |    Protocol    |                      |
     +-------------+    Stack(s)    +----------------------+
                   +-------+--------+
                           V
   +-----------------------------------------------------+
   |               Network-Layer Interface               |
   +-----------------------------------------------------+
    \end{lstlisting}
    \caption{TAPS architecture as described by RFC 9621 \cite{paulyArchitectureRequirementsTransport2025}}
\end{figure}

\section{Development Timeline} \label{projectTimeline}
The first goal while developing the CTaps library was to get a minimal viable implementation,
able to send and receive UDP messages to a ping server. In order to
achieve this, the project structure was generated to allow for automated
testing in order to verify this behaviour as easily as possible.
A simple dummy library was developed, both to test the packaging system,
and to verify the testing library Google Test and the mocking framework FFF.
After both libraries and the project structure were verified to work
as expected, minimal versions of the Preconnection, Connection and Message
were defined in order to implement the functionality needed to set up a
basic UDP ping test.

After the basic structures needed to send an UDP message were implemented,
a test verifying this behaviour was implemented. Libuv was also selected
as the library to use for the primary event loop. This would serve as the
basis for CTaps' asynchronous, event-driven model. Libuv had the added benefit
of having built-in support for UDP and TCP handles. This eased the development
of the initial application pinging a UDP ping server, which was complete around
a week after starting development. A simple interface for future protocols was
also implemented, with members  pointers for each function necessary. This would
make it as easy as possible to implement a new transport layer protocol, without
changing the actual interface of the library or adding new logic outside of the
new protocol implementation.

\begin{figure}[H]
\begin{minted}{C}
typedef struct ProtocolImplementation {
  const char* name;
  ProtocolFeatures features;
  int (*init)(struct Connection* connection, InitDoneCb init_done_cb);
  int (*send)(struct Connection*, Message*);
  int (*receive)(struct Connection*, ReceiveMessageRequest receive_cb);
  int (*listen)(struct SocketManager* socket_manager);
  int (*stop_listen)(struct SocketManager*);
  int (*close)(const struct Connection*);
} ProtocolImplementation;
\end{minted}
\label{The interface used to interact with protocols}
\end{figure}

After writing a successful test both sending
and receiving a message from an UDP server, development started on simple DNS
lookup functionality, using the built-in asynchronous implementations of
\texttt{getaddrinfo} from Libuv.

After the development of the basic DNS functionality was completed, 
the listening functionality was started, in order to try to replace
the external UDP ping server initially used in the tests. The listening
functionality necessitated a fair amount of redesign and new features.
For example, there was previously no support for incoming Messages being
multiplexed to a Connection, as Connections were bound to a single
Libuv UDP handle each. In accordance with the methodology described in
section \ref{developmentMethodology}, the minimal changes needed to
make Connections abstract enough to support more than one protocol,
and the needed code to multiplex incoming UDP messages were implemented.
A new SocketManager structure was introduced, in order to handle listening
for connections. It was designed to multiplex the incoming UDP messages
and delegate incoming TCP connections, for when TCP functionality would
be implemented.

After the necessary updates to Connections were made in order
for Listeners to work, work was started on more advanced candidate
gathering and candidate racing algorithms. The previous algorithm
only considered which selection properties were prohibited/required,
ranked protocols according to the number of avoid/prefer categories
it fulfilled, picking the top ranked protocol every time. A more advanced
Candidate Gathering algorithm with lay the groundwork for later Candidate
Racing, ensuring that only branches which are actually prohibited are
eliminated, making the overal solution far more flexible and robust,
allowing the use of a undesirable path, if none of the more desirable
ones succeed. The Candidate Gathering functionality was implemented using
N-ary tree functionality from Glib. This was chosen because the robustness
of an established library, while the Candidate Gathering functionality would
not be directly exposed to consuming application, making sure that the
actual API was kept neat and in pure C.

The Candidate Racing algorithm was implemented according to the recommendations
in RFC 9623. Cont...
